\begin{abstract}
The discovery of advanced materials is fundamentally limited by the slow, trial-and-error nature of experimental synthesis. While Large Language Models (LLMs) promise to accelerate this process by integrating vast domain knowledge, we reveal a critical failure mode known as \textbf{\textit{contextual tunneling}}, wherein naive knowledge integration causes LLMs to over-anchor on narrow retrieval paths---such as irrelevant synthesis recipes---while suppressing broader physical principles. Through a systematic evaluation on \textbf{dopant control in 2D materials}, we demonstrate that naive knowledge graph augmentation degrades performance by 20--35\% on inverse design tasks compared to direct prompting.

To address this challenge, we introduce \texttt{ARIA} (Autonomous Reasoning Intelligence for Atomics), a causal-aware framework tailored for materials discovery. \texttt{ARIA} features: (i) hierarchical reasoning that leverages the Processing-Structure-Property (PSP) paradigm to handle knowledge sparsity; (ii) enhanced analogical transfer that adapts synthesis protocols using physicochemical similarity metrics; and (iii) dynamic knowledge enrichment through online searching. Extensive experiments show that while naive KG integration consistently underperforms baseline LLMs in synthesis planning, \texttt{ARIA} not only recovers this loss but also provides verifiable causal pathways. This transforms the "black box" suggestions of standard LLMs into the interpretable, physically grounded experimental protocols required for trustworthy materials design.
\end{abstract}